% Options for packages loaded elsewhere
% Options for packages loaded elsewhere
\PassOptionsToPackage{unicode}{hyperref}
\PassOptionsToPackage{hyphens}{url}
\PassOptionsToPackage{dvipsnames,svgnames,x11names}{xcolor}
%
\documentclass[
  spanish,
  12pt,
  a4paper,
]{article}
\usepackage{xcolor}
\usepackage[top=1in,bottom=1in,left=1in,right=1in]{geometry}
\usepackage{amsmath,amssymb}
\setcounter{secnumdepth}{-\maxdimen} % remove section numbering
\usepackage{iftex}
\ifPDFTeX
  \usepackage[T1]{fontenc}
  \usepackage[utf8]{inputenc}
  \usepackage{textcomp} % provide euro and other symbols
\else % if luatex or xetex
  \usepackage{unicode-math} % this also loads fontspec
  \defaultfontfeatures{Scale=MatchLowercase}
  \defaultfontfeatures[\rmfamily]{Ligatures=TeX,Scale=1}
\fi
\usepackage{lmodern}
\ifPDFTeX\else
  % xetex/luatex font selection
\fi
% Use upquote if available, for straight quotes in verbatim environments
\IfFileExists{upquote.sty}{\usepackage{upquote}}{}
\IfFileExists{microtype.sty}{% use microtype if available
  \usepackage[]{microtype}
  \UseMicrotypeSet[protrusion]{basicmath} % disable protrusion for tt fonts
}{}
% Make \paragraph and \subparagraph free-standing
\makeatletter
\ifx\paragraph\undefined\else
  \let\oldparagraph\paragraph
  \renewcommand{\paragraph}{
    \@ifstar
      \xxxParagraphStar
      \xxxParagraphNoStar
  }
  \newcommand{\xxxParagraphStar}[1]{\oldparagraph*{#1}\mbox{}}
  \newcommand{\xxxParagraphNoStar}[1]{\oldparagraph{#1}\mbox{}}
\fi
\ifx\subparagraph\undefined\else
  \let\oldsubparagraph\subparagraph
  \renewcommand{\subparagraph}{
    \@ifstar
      \xxxSubParagraphStar
      \xxxSubParagraphNoStar
  }
  \newcommand{\xxxSubParagraphStar}[1]{\oldsubparagraph*{#1}\mbox{}}
  \newcommand{\xxxSubParagraphNoStar}[1]{\oldsubparagraph{#1}\mbox{}}
\fi
\makeatother


\usepackage{longtable,booktabs,array}
\newcounter{none} % for unnumbered tables
\usepackage{calc} % for calculating minipage widths
% Correct order of tables after \paragraph or \subparagraph
\usepackage{etoolbox}
\makeatletter
\patchcmd\longtable{\par}{\if@noskipsec\mbox{}\fi\par}{}{}
\makeatother
% Allow footnotes in longtable head/foot
\IfFileExists{footnotehyper.sty}{\usepackage{footnotehyper}}{\usepackage{footnote}}
\makesavenoteenv{longtable}
\usepackage{graphicx}
\makeatletter
\newsavebox\pandoc@box
\newcommand*\pandocbounded[1]{% scales image to fit in text height/width
  \sbox\pandoc@box{#1}%
  \Gscale@div\@tempa{\textheight}{\dimexpr\ht\pandoc@box+\dp\pandoc@box\relax}%
  \Gscale@div\@tempb{\linewidth}{\wd\pandoc@box}%
  \ifdim\@tempb\p@<\@tempa\p@\let\@tempa\@tempb\fi% select the smaller of both
  \ifdim\@tempa\p@<\p@\scalebox{\@tempa}{\usebox\pandoc@box}%
  \else\usebox{\pandoc@box}%
  \fi%
}
% Set default figure placement to htbp
\def\fps@figure{htbp}
\makeatother



\ifLuaTeX
\usepackage[bidi=basic,shorthands=off]{babel}
\else
\usepackage[bidi=default,shorthands=off]{babel}
\fi
\ifLuaTeX
  \usepackage{selnolig} % disable illegal ligatures
\fi


\setlength{\emergencystretch}{3em} % prevent overfull lines

\providecommand{\tightlist}{%
  \setlength{\itemsep}{0pt}\setlength{\parskip}{0pt}}





%-----------------------------------------------------------------------------%
%                            SOME PREAMBLE STUFF                              %
%-----------------------------------------------------------------------------%

%-----------------------------------------------------------------------------%
%                            LOAD PACKAGES                                    %
%-----------------------------------------------------------------------------%

\usepackage{hyperref}
\usepackage{setspace}
\usepackage{geometry}
\usepackage[style=apa, backend=biber,annotation=false]{biblatex}
\usepackage[x11names]{xcolor}
\usepackage{longtable}
\usepackage{fancyhdr} % for header
\usepackage{soul} % to make highlight for tags
\usepackage{pdfpages} % to append pdf pages
\usepackage{fontspec} % for fonts

%-----------------------------------------------------------------------------%
%                            Header                                           %
%-----------------------------------------------------------------------------%

\setlength{\headheight}{38pt}
\setmainfont{Helvetica Neue}
\makeatletter
\@ifpackageloaded{tcolorbox}{}{\usepackage[skins,breakable]{tcolorbox}}
\@ifpackageloaded{fontawesome5}{}{\usepackage{fontawesome5}}
\definecolor{quarto-callout-color}{HTML}{909090}
\definecolor{quarto-callout-note-color}{HTML}{0758E5}
\definecolor{quarto-callout-important-color}{HTML}{CC1914}
\definecolor{quarto-callout-warning-color}{HTML}{EB9113}
\definecolor{quarto-callout-tip-color}{HTML}{00A047}
\definecolor{quarto-callout-caution-color}{HTML}{FC5300}
\definecolor{quarto-callout-color-frame}{HTML}{acacac}
\definecolor{quarto-callout-note-color-frame}{HTML}{4582ec}
\definecolor{quarto-callout-important-color-frame}{HTML}{d9534f}
\definecolor{quarto-callout-warning-color-frame}{HTML}{f0ad4e}
\definecolor{quarto-callout-tip-color-frame}{HTML}{02b875}
\definecolor{quarto-callout-caution-color-frame}{HTML}{fd7e14}
\makeatother
\makeatletter
\@ifpackageloaded{caption}{}{\usepackage{caption}}
\AtBeginDocument{%
\ifdefined\contentsname
  \renewcommand*\contentsname{Tabla de contenidos}
\else
  \newcommand\contentsname{Tabla de contenidos}
\fi
\ifdefined\listfigurename
  \renewcommand*\listfigurename{Listado de Figuras}
\else
  \newcommand\listfigurename{Listado de Figuras}
\fi
\ifdefined\listtablename
  \renewcommand*\listtablename{Listado de Tablas}
\else
  \newcommand\listtablename{Listado de Tablas}
\fi
\ifdefined\figurename
  \renewcommand*\figurename{Figura}
\else
  \newcommand\figurename{Figura}
\fi
\ifdefined\tablename
  \renewcommand*\tablename{Tabla}
\else
  \newcommand\tablename{Tabla}
\fi
}
\@ifpackageloaded{float}{}{\usepackage{float}}
\floatstyle{ruled}
\@ifundefined{c@chapter}{\newfloat{codelisting}{h}{lop}}{\newfloat{codelisting}{h}{lop}[chapter]}
\floatname{codelisting}{Listado}
\newcommand*\listoflistings{\listof{codelisting}{Listado de Listados}}
\makeatother
\makeatletter
\makeatother
\makeatletter
\@ifpackageloaded{caption}{}{\usepackage{caption}}
\@ifpackageloaded{subcaption}{}{\usepackage{subcaption}}
\makeatother
\usepackage{bookmark}
\IfFileExists{xurl.sty}{\usepackage{xurl}}{} % add URL line breaks if available
\urlstyle{same}
\hypersetup{
  pdftitle={Tinción de gotas de lípidos en adipocitos con rojo oleoso},
  pdfauthor={Sara y Dianis},
  pdflang={es-MX},
  colorlinks=true,
  linkcolor={DarkSlateBlue},
  filecolor={Maroon},
  citecolor={DarkSlateBlue},
  urlcolor={DarkSlateBlue},
  pdfcreator={LaTeX via pandoc}}


%-----------------------------------------------------------------------------%
%                            BIBLATEX                                         %
%-----------------------------------------------------------------------------%

\usepackage[style=apa,backend=biber]{biblatex}

\setlength\bibitemsep{0pt}  % No space between bib entries
\renewcommand*{\bibfont}{\footnotesize}  % Use smaller font
\setlength\bibhang{\parindent}  % Match document indentation

% Fix biblatex's odd preference for using In: by default.
\renewbibmacro{in:}{%
  \ifentrytype{article}{}{%
  \printtext{\bibstring{}\intitlepunct}}}

\addbibresource{refs.bib}

\begin{document}
%-----------------------------------------------------------------------------%
%                            SOME PREAMBLE STUFF                              %
%-----------------------------------------------------------------------------%

%-----------------------------------------------------------------------------%
%                            Set Page STYLE                                   %
%-----------------------------------------------------------------------------%

\pagestyle{fancy}
\fancyhead[L]{\thepage \\ Sara y Dianis }
\fancyhead[R]{Tinción de gotas de lípidos en adipocitos con rojo
oleoso \\ Laboratorio de Inmunometabolismo }
\cfoot{\textbf{Version} \texttt{Inicial} 2022-02 }
\renewcommand*\contentsname{Tabla de contenidos}
{
\setcounter{tocdepth}{3}
\tableofcontents
}

\newpage{}

\subsubsection{Descripción del
Protocolo}\label{descripciuxf3n-del-protocolo}

\begin{tcolorbox}[enhanced jigsaw, arc=.35mm, bottomrule=.15mm, bottomtitle=1mm, breakable, colback=white, colbacktitle=quarto-callout-note-color!10!white, colframe=quarto-callout-note-color-frame, coltitle=black, left=2mm, leftrule=.75mm, opacityback=0, opacitybacktitle=0.6, rightrule=.15mm, title=\textcolor{quarto-callout-note-color}{\faInfo}\hspace{0.5em}{Pasos}, titlerule=0mm, toprule=.15mm, toptitle=1mm]

\begin{itemize}
\tightlist
\item
  Retirar el medio a las células.
\item
  Agregar por las paredes del pozo 1 mL de PBS 1X estéril a las células,
  realizar movimientos circulares y tenues, para eliminar el medio.
  Repetir este paso 2 veces más.
\item
  Retirar el PBS 1X y agregar 700 μL de paraformaldehído (PAF) al 4\%,
  incubar a temperatura ambiente por 20 min.
\item
  Eliminar el PAF con la pipeta y realizar 3 lavados con PBS 1X, ya
  descritos anteriormente.
\item
  Agregar 500 μL de rojo oleoso a la monocapa e incubar a temperatura
  ambiente por 1 hora.
\item
  Retirar el rojo oleoso y realizar 3 lavados con PBS 1X.
\item
  Observar al microscopio.
\end{itemize}

\end{tcolorbox}

\begin{tcolorbox}[enhanced jigsaw, arc=.35mm, bottomrule=.15mm, bottomtitle=1mm, breakable, colback=white, colbacktitle=quarto-callout-caution-color!10!white, colframe=quarto-callout-caution-color-frame, coltitle=black, left=2mm, leftrule=.75mm, opacityback=0, opacitybacktitle=0.6, rightrule=.15mm, title=\textcolor{quarto-callout-caution-color}{\faFire}\hspace{0.5em}{Notas}, titlerule=0mm, toprule=.15mm, toptitle=1mm]

\begin{itemize}
\tightlist
\item
  La adición y eliminación de los reactivos se realizará por las paredes
  del pozo, para evitar el desprendimiento de la monocapa de células.
\item
  Para los lavados, los movimientos no deberán ser vigorosos porque
  puede provocar el desprendimiento de la monocapa de células.
\item
  Los volúmenes que se mencionan en el procedimiento es para el tamaño
  de un pozo de placa 6 pozos, estos volúmenes se pueden adecuar de
  acuerdo al tamaño del pozo en donde se esté realizando la tinción.
\end{itemize}

\end{tcolorbox}

\subsubsection{Diagrama de Flujo}\label{diagrama-de-flujo}

\begin{figure}[H]

{\centering \includegraphics[width=0.25\linewidth,height=\textheight,keepaspectratio]{images/oil.png}

}

\caption{Figura 1. Diagrama de flujo}

\end{figure}%

\subsubsection{Anexos}\label{anexos}

Soluciones

Preparación de medios. a) Medio de crecimiento

Medio DMEM + 10\% SFB + antibiótico/antimicótico + l-glutamina

\begin{enumerate}
\def\labelenumi{\alph{enumi})}
\setcounter{enumi}{1}
\item
  Medios de diferenciación(MD) La siguiente tabla es descrita para 3 Ml
  de m, es decir para un pozo. Al medio de crecimiento se le adicionara
  los reactivos necesarios para preparar los medios de diferenciación

  MD-I MD-II MD-III Dexamentasona 25 μM 30 μL 30 μL --- IBMX 125 μM 12
  μL 12 μL --- Insulina 300 UI/mL 1 μL 1 μL 1 μL Rosiglitazona 20 mM ---
  3 μL de una dilución 1:10 (9 μL de medio de crecimiento + 1 μL de
  rosiglitazona 20 Mm) ---
\end{enumerate}

\subsubsection{\texorpdfstring{Documentos de
\textbf{?@sec-ref}}{Documentos de ?@sec-ref}}\label{documentos-de--sec-ref}

Zebisch K, Voigt V, Wabitsch M, Brandsch M. Protocol for effective
differentiation of 3T3-L1 cells to adipocytes. Analytical biochemistry.
2012;425(1):88-90.


\printbibliography[title=References]



\end{document}
